%% Needed for the postscript graphics.
%%\input epsf

%% My macros.
\input ../util/macro

%% The length of a page is 10in.
\vsize=9.5in
\voffset=-.3in

%\font\helv = phvr
%\font\times = ptmr
%\font\zap = pzcmi at 12pt
%\font\pal = pplro at 12pt

%% Don't indent paragraphs.
\parindent=0pt
\parskip=\baselineskip

\mantitle{The Derivative Revised}
{R. Scott McIntire} {Mar 01 2004}

\section{Another Look at the Derivative}
We have seen that the derivative at any point of a functions is defined 
to be the slope of the 
tangent line to the graph of the function. If we know the derivative at 
a point we can approximate the value of a function near that point.
For $h$ small we have
$$
f'(x) \approx {f(x+h) - f(x) \over h }
$$
Or, after multiplying by $h$ and rearranging we may write
$$
f(x+h) \approx f(x) + h f'(x)
$$
We can replace the $\approx$ sign with and equal sign if we specify 
the error term. 
$$
f(x+h) = f(x) + h f'(x) + {\rm error\/}(h)
$$
Notice that this error function has the property that 
$\lim_{h \rightarrow 0} {{\rm error\/}(h) \over h } = 0$ 
The graph below illustrates the approximation. 
% TBD graph of deriv approximation.

From this we see that the tangent line is an approximation to the function $f$.
I think you can see from the picture that there are other linear approximations 
but the tangent line is the best in the sense that the error left over is 
``smaller'' than linear. By this we mean 
$\lim_{h \rightarrow 0} {{\rm error\/}(h) \over h} = 0$. 
We can think of our approximation to $f(x+h)$ as $f(x)$ plus a linear function 
of $h$ which represents a tangent line through the origin, $f'(x) h$.
That is, the linear function at $x$ through the origin we will call $DF(x)$, 
defined as
$$
DF(x)(h) = f'(x) h
$$ 
This is just a restatement of the idea that a straight line through the 
origin is the graph of a linear function $l(x) = m x$. The $m$ in this 
case is the slope of $f$ at $x$, $f'(x)$.

Here is a better version of the definition of the derivative of a function, $f$.
It gets us back to the geometry of the situation. 
\proclaim Definition(Derivative). The derivative of a function $f:R \mapsto R$ 
at $x$, denoted $Df(x)$, is a linear function from $R$ to $R$ 
that satisfies the equation
$$
f(x+h) = f(x) + Df(x)(h) + {\rm error\/}(h)
$$
where error satisfies 
$\lim_{h \rightarrow 0} {{\rm error\/}(h) \over h} = 0$\footnote{\kern 0.5pt \raise 0.5ex \hbox{\dag}}
{As example of such an error function is ${\rm error\/}(h) = h^2$.}.

This definition says that the derivative of a function at a point is 
the best linear approximation to the function. 

What does this say about our maximization problem. Recall we needed a 
measure of steepness that was computable and 0 when the tangent line was flat.
What we need is a way to recognize when the tangent line is flat. With 
our new notion of derivative, the tangent line is flat when the derivative 
is the zero function. Since the zero linear function corresponds to the 
number 0 (the slope of the line); and, more generally, each 
linear function corresponds to the slope of the tangent line we have separated 
the notion of a derivative into a geometric object and then have 
a representation of that object that is computable. We formalize this 
representation as:

\proclaim Theorem(Representation of the linear function $L:R \mapsto R$). 
Each linear function
$L:R \mapsto R$ has the from $L(x) = m x$. 

We abuse terminology and will refer to linear function, $Df(x)$, and number 
$f'(x)$ as the derivative of $f$ at $x$. 

What is the point of all of this. This doesn't improve the situation with 
our example problem; it only makes things more complex. The reason this 
is a better definition is that we see the separation of the geometry and 
the representation of such in terms of algebra. This is important when 
we try to generalize the notion of a derivative to functions other than 
$f:R \mapsto R$ where we are unable to visualize the geometry. 
We now state a more general definition of a derivative.

\proclaim Definition(Derivative). The derivative of a function 
${\bf f\/}:R^n \mapsto R^m$ 
at $x$, denoted $D{\bf f\/}({\bf x\/})$, is a linear function 
from $R^n$ to $R^m$ 
that satisfies the equation
$$
{\bf f\/}({\bf x+h\/}) ={\bf f\/}({\bf x\/}) + 
D{\bf f\/}({\bf x\/})({\bf h\/}) + {\rm error\/}({\bf h\/})
$$
where the 'error' function satisfies 
$\lim_{\|{\bf h\/}\| \rightarrow 0} 
{\|{\rm error\/}({\bf h\/})\| \over \|{\bf h\/}\|} = 0$\footnote{\kern 0.5pt \raise 0.5ex \hbox{\ddag}}
{As example of such an error function is ${\rm error\/}(h) = h^2$.}.


We list some representation theorems that provide a computational object for 
the derivative.

\proclaim Theorem(Representation of the linear function L:$R^n\mapsto R$). 
Each linear function
$L:R^n \mapsto R$ is represented by a vector in $R^n$ by
 $L({\bf x\/}) = {\bf m\/} {\bf \cdot\/} {\bf x\/}$. 

\proclaim Definition(Gradient). The vector that represents the derivative 
of $f: R^n \mapsto R$ at ${\bf x\/}$ is called the gradient of $f$ 
at ${\bf x\/}$ and is denoted, 
$\nabla f({\bf x\/})$. If Cartesian coordinates are used then the components 
of the gradient are 
$$
\nabla f({\bf x\/}) = \left[{\partial f({\bf x\/}) \over \partial x_1}, 
{\partial f({\bf x\/}) \over \partial x_2}, \ldots, 
{\partial f({\bf x\/}) \over \partial x_n}\right]^T 
$$

\proclaim Definition(Jacobian). The matrix that represents the derivative 
of ${\bf f\/}: R^n \mapsto R^m$ at ${\bf x\/}$ is called the 
Jacobian matrix of ${\bf f\/}$ 
at ${\bf x\/}$ and is denoted, 
$J {\bf f\/}({\bf x\/})$. If Cartesian coordinates are
used then the $i^{\rm th\/}$
$j^{\rm th\/}$ component is
$$
\left[J{\bf f\/}({\bf x\/})\right]_{ij} = {\partial f_i \over \partial x_j}
$$



\bye

